\documentclass[12pt,a4paper]{article}
\usepackage[margin=2.3cm]{geometry}
\usepackage{fontspec}
\setmainfont{Times New Roman}
\usepackage{setspace}
\singlespacing
\usepackage{titlesec}
\titleformat{\section}{\normalfont\large\bfseries}{}{0em}{}
\titleformat{\subsection}{\normalfont\normalsize\bfseries}{}{0em}{}
\titlespacing{\section}{0pt}{10pt}{4pt}
\titlespacing{\subsection}{0pt}{8pt}{2pt}
\usepackage{hyperref}
\hypersetup{hidelinks}
\setlength{\parindent}{0pt}
\setlength{\parskip}{4pt}
\newenvironment{refs}{\begin{list}{}{\setlength{\leftmargin}{1.5em}\setlength{\itemindent}{-1.5em}\setlength{\itemsep}{1pt}}}{\end{list}}

\begin{document}

{\large\bfseries BE451 Degree Project in Finance --- Thesis Synopsis, Fall 2026}

\textbf{Students:} [Full name (student ID)], [Full name (student ID)] \\
\textbf{Preferred tutor:} [Michael Klug / Alvin Chen / no preference]

\vspace{4pt}\hrule\vspace{4pt}

\section{Topic 1: Do Taxes Depress or Stimulate Corporate Borrowing? Evidence from Sweden's 2019 Tax Cut and Interest Deduction Limit}

\subsection{(i) Tentative research question}

\emph{Did Swedish private firms reduce or increase leverage after the 2019 reform that cut the corporate tax rate and capped interest deductibility, and does the response differ between firms that only received a lower rate and firms that lost part of their tax shield?}

The reform (prop.\ 2017/18:245, in force 1 January 2019) contains two treatments with opposite predictions under the trade-off theory of capital structure. The statutory rate fell from 22 to 21.4 percent in 2019 and to 20.6 percent in 2021, lowering the value of the interest tax shield for every firm. Simultaneously, net interest expense became deductible only up to 30 percent of taxable EBITDA, with a simplification rule that allows SEK 5 million of net interest per group of associated companies regardless. Trade-off theory (Heider and Ljungqvist, 2015) predicts lower leverage from both changes. Ivanov, Pettit, and Whited (2026) find the opposite for small US private firms: after tax cuts, higher after-tax cash flows move default further away, credit spreads fall, and leverage rises. Sanati and Beyhaghi (2025) show that losing the deduction lowers debt and investment without equity substitution. Sweden lets us run both tests in a single population of all limited companies: the cap binds only for firms whose net interest exceeds both SEK 5 million and 30 percent of EBITDA, which creates a kink in the marginal tax treatment of interest at SEK 5 million and a clear split between bound and unbound firms, while the rate cut is a common shock whose bite varies with pre-reform taxable profits. A null result on both margins would itself be informative: Swedish private-firm capital structure would then not be tax-driven at the elasticities estimated in the US literature.

\subsection{(ii) Potential data sources and empirical methods}

\textbf{Data.} Serrano (Swedish House of Finance), annual accounts for all Swedish limited companies 2013--2024: total and interest-bearing debt, interest expense, tax expense, EBITDA, investment, dividends, group parent and intra-group balances. Access through SWAMID [verified on (date); N $\approx$ (number) firm-years]. Treatment variables: (a) rate cut: pre-reform effective tax payments over assets, with loss-making firms untreated; (b) cap: net interest expense in 2017--2018 above versus below SEK 5 million, aggregated within groups of associated companies from Serrano's group identifiers; (c) rental-housing firms, which received a 12 percent primary deduction in the same reform, as a third treatment. The cap binds for a few thousand firms, mainly in real estate and private-equity-owned groups.

\textbf{Methods.} For the rate cut, difference-in-differences of leverage, new borrowing, interest expense over debt and investment on pre-reform tax exposure interacted with post-2019 and post-2021 indicators, with firm and industry-by-year fixed effects. For the cap, a bunching analysis at the SEK 5 million kink in group net interest (the marginal deduction falls from full to zero at the kink for firms whose EBITDA-based allowance is below SEK 5 million), and difference-in-differences for firms bound by the cap in 2017--2018 versus unbound firms of similar size and leverage. Identification rests on the law fixing both the rate path and the kink: within industry and year, firms with different pre-reform exposure faced the same macroeconomic conditions. Threats and mitigations: manipulation around the threshold (bunching is itself a finding; treatment assigned on pre-reform net interest); intra-group reallocation of interest (group-level aggregation, following Brok); the 2022 rate cycle (industry-by-year effects and a 2016--2018 placebo); concurrent 2019 hybrid and CFC rules (exclude large multinationals). Benchmark: the Heider--Ljungqvist elasticity implies a leverage change of about 0.6 percentage points for the 1.4 point rate cut; with population data the minimum detectable effect is well below that. Three-month plan: the rate-cut analysis first; then the bunching analysis, the bound-versus-unbound comparison and the group channel. Fallback: rate-cut analysis only.

\subsection{(iii) Short literature review}

\textbf{Ivanov, Pettit, and Whited (2026, Review of Economic Studies)} use US state tax changes and private-firm data and find that leverage rises after tax cuts because lower taxes raise the productivity of capital and lower credit spreads. We replicate their tax-cut test with the 2019 and 2021 rate steps and extend it by adding the simultaneous deduction cap, which their setting lacks. \textbf{Sanati and Beyhaghi (2025, Review of Financial Studies)} exploit the USD 25 million exemption of the 2017 US reform in a regression discontinuity and find that losing the tax benefit of debt lowers debt and investment. We adapt their exemption-threshold design to Sweden's simplification rule, which is a kink rather than a discontinuity, using bunching at the kink and a bound-versus-unbound comparison at the group level. \textbf{Heider and Ljungqvist (2015, Journal of Financial Economics)} estimate from US state tax changes that a one percentage point tax increase raises leverage by about 0.4 percentage points and that responses are asymmetric; we use their estimate as the benchmark for our minimum detectable effect. \textbf{Brok (2026, Journal of Financial and Quantitative Analysis)} shows that tax avoidance acts as a non-debt tax shield that crowds out the debt tax shield, so the leverage response to taxes is weaker for high-avoidance firms; \textbf{Brok (2024, Review of Finance)} shows that corporate group structure creates additional tax incentives for leverage because parents can deduct subsidiary losses. Both imply that tax effects on debt must be measured at the group level, which Serrano's group identifiers allow, and let us test whether groups reallocate interest across entities to stay under the cap. Relative to these papers, we replicate both the tax-cut and the deduction-cap designs in one reform and one population, allowing us to test the trade-off and productivity channels against each other with a threshold that binds at the group level. Our contribution is threefold: the first evidence on Sweden's 2019 reform for the full private-firm population; a horse race between two channels with opposite signs in the same data; and a hand-built group-level net-interest variable that assigns the kink correctly, which firm-level data cannot.

\subsection*{References}
\begin{refs}
\item Brok, Peter, 2024, Debt and taxes: The role of corporate group structure, \emph{Review of Finance} 28, 2051--2082.
\item Brok, Peter, 2026, Debt and taxes: The role of tax avoidance, \emph{Journal of Financial and Quantitative Analysis} 61, 2007--2032.
\item Heider, Florian, and Alexander Ljungqvist, 2015, As certain as debt and taxes: Estimating the tax sensitivity of leverage from state tax changes, \emph{Journal of Financial Economics} 118, 684--712.
\item Ivanov, Ivan T., Luke Pettit, and Toni M. Whited, 2026, Taxes depress corporate borrowing: Evidence from private firms, \emph{Review of Economic Studies} 93, 2715--2749.
\item Sanati, Ali, and Mehdi Beyhaghi, 2025, How does removing the tax benefits of debt affect firms? Evidence from the 2017 U.S. tax reform, \emph{Review of Financial Studies}, forthcoming (advance access, doi:10.1093/rfs/hhaf041).
\end{refs}

\newpage

\section{Topic 2: Forced Fund Flows and Price Pressure in Swedish Equities: Evidence from the Premium Pension Platform Procurement}

\subsection{(i) Tentative research question}

\emph{When the state moves tens of billions of kronor between mutual funds on a pre-announced administrative timetable, do the stocks those funds must sell and buy experience price pressure and subsequent reversal, or does the market absorb the flows---and is it net demand (after netting sellers against buyers) or gross turnover that matters?}

Since 2024 the Swedish Fund Selection Agency has procured the premium pension fund platform category by category. Funds that lose the procurement are deregistered and savers' capital is transferred automatically to procured funds unless the saver chooses otherwise; the Pensions Agency executes the sales as ``split fund closures'', typically one tranche per week over weeks or months. The Swedish active-equity category (awarded 27 August 2025; SEK 92 billion; 573,000 savers; about 25 losing funds closed between 10 November 2025 and 21 April 2026, roughly one fund per week) and the Nordic small-cap category (2025) are complete. The Swedish small-cap category (awarded 28 May 2026; SEK 34 billion; about 180,000 savers) is being executed now: nine losing funds holding about SEK 8.3 billion became non-selectable on 18 August 2026 and their capital flows mainly to five newly procured funds. These flows are large relative to the liquidity of Swedish small and mid caps, their timing and size are set administratively rather than by the stocks' fundamentals, and both the selling and the receiving portfolios are observable. The setting therefore lets us test the price-pressure hypothesis of Coval and Stafford (2007) and Lou (2012) with an exogenous flow whose netting across buyers and sellers we can measure. A null is informative: pre-announced, staggered institutional flows would then be absorbed without price effects, which speaks directly to the debate on how elastic demand for individual stocks is (Koijen and Yogo, 2019).

\subsection{(ii) Potential data sources and empirical methods}

\textbf{Data.} Fund holdings by ISIN from Finansinspektionen's quarterly ``Fondinnehav'' files (public, 2018Q4--2026Q2), complemented by funds' semi-annual reports and Morningstar Direct through the Swedish House of Finance for foreign-domiciled funds. Fund-level premium pension capital and monthly net trading per fund from the Pensions Agency's monthly statistics (public), and deregistration dates from its fund list. Daily prices, volumes and, where available, bid--ask spreads from FinBas and LSEG Workspace. Tranche timing reconstructed from daily net flows of the losing funds' premium pension share classes. Sample: all completed waves with Swedish equity exposure (Swedish active, Swedish passive, Nordic active, Nordic small cap; 2024--2026) as the main sample, and the Swedish small-cap wave (August--December 2026) as an out-of-sample test using second-quarter holdings. Expected sample: several hundred fund--stock--tranche observations across roughly 200--300 stocks [verify data access on (date)].

\textbf{Methods.} For each stock and tranche we compute gross flow (holdings of closing funds times capital to be sold) and net demand (gross sales minus purchases implied by the receiving funds' portfolios), scaled by average daily volume, following the corrections in Wardlaw (2020). We estimate an event study of abnormal returns, volume and spreads over days $-20$ to $+30$ around each tranche, and panel regressions with stock and day fixed effects of abnormal returns on net demand and of volume and spreads on gross flow, with interactions for liquidity and for the overlap between losing and receiving portfolios. Controls: Swedish small and mid caps without premium pension exposure, matched on size and liquidity. Identification: the timing and size of each flow are set by the procurement schedule and the Pensions Agency's execution rule, not by the stock; the staggered order across weeks gives within-stock variation. Threats and mitigations: non-random selection of winners (irrelevant for a timing-based test; tilts controlled with matched non-exposed stocks and a placebo in the announcement quarter before any flow); pre-trading by losing funds (attenuates toward zero; tested with pre-event drift and the change in holdings between quarters); overlap netting out demand (measured explicitly, and it is the mechanism test); data lag for the live wave (main results on completed waves). Three-month plan: the Swedish active wave first (holdings, flows and event study), then the pooled waves and the small-cap wave as it unfolds. Fallback: completed waves only.

\subsection{(iii) Short literature review}

\textbf{Coval and Stafford (2007, Journal of Financial Economics)} show that mutual funds facing large outflows sell holdings at temporarily depressed prices that reverse within months; we replicate the fire-sale event study with flows that are administratively forced rather than driven by performance, which removes the reverse-causality concern in their setting. \textbf{Lou (2012, Review of Financial Studies)} builds a flow-induced trading measure from fund holdings and shows it predicts returns and reversals; we replicate his measure and extend it by netting the buying side, which is unobservable in his data but known here. \textbf{Wardlaw (2020, Journal of Finance)} shows that standard flow-pressure measures mechanically embed returns and proposes a corrected measure; we adopt his correction. \textbf{Dahlquist, Martinez, and Söderlind (2017, Review of Financial Studies)} study Swedish premium pension investors' activity and performance and document the institutional setting; we use their description of the platform and its savers. \textbf{Koijen and Yogo (2019, Journal of Political Economy)} estimate a demand system in which institutional demand for individual stocks is inelastic; our flows are a direct test of that elasticity outside the United States. Relative to Coval and Stafford (2007) and Lou (2012), we replicate the price-pressure design on flows whose timing, size and counterparties are set by an administrative procurement, allowing us to separate net demand from gross turnover and to test price pressure against absorption. Our contribution is threefold: the first evidence on price effects of Sweden's premium pension platform procurement; a flow measure that nets sellers against known buyers; and a live out-of-sample test on the small-cap wave executed during the thesis period. Benchmark: Coval and Stafford report abnormal returns of several percent around forced sales with reversal; with hundreds of stock--tranche observations our minimum detectable effect is a fraction of that.

\subsection*{References}
\begin{refs}
\item Coval, Joshua, and Erik Stafford, 2007, Asset fire sales (and purchases) in equity markets, \emph{Journal of Financial Economics} 86, 479--512.
\item Dahlquist, Magnus, Jos\'e Vicente Martinez, and Paul S\"oderlind, 2017, Individual investor activity and performance, \emph{Review of Financial Studies} 30, 866--899.
\item Koijen, Ralph S. J., and Motohiro Yogo, 2019, A demand system approach to asset pricing, \emph{Journal of Political Economy} 127, 1475--1515.
\item Lou, Dong, 2012, A flow-based explanation for return predictability, \emph{Review of Financial Studies} 25, 3457--3489.
\item Wardlaw, Malcolm, 2020, Measuring mutual fund flow pressure as shock to stock returns, \emph{Journal of Finance} 75, 3221--3243.
\end{refs}

\end{document}
